\documentclass{article}
\usepackage{graphicx} % Required for inserting images
\usepackage{amsmath}
\usepackage{amssymb}
\title{Differential Geometry}
\author{Maxwel Kim}
\date{June 2026}

\begin{document}

\maketitle

\section*{curvature}

\textbf{Velocity vector} 

\[
T(s)=\gamma'(s)
\]

has unit length and is tangent to the curve at point

\[
p=\gamma(s)
\]

and curvature at p is

\[
T'(s)=\gamma''(s)
\]

therefore.

\textbf{signed curvature or simply curvature} is 

\[
\kappa =<T',n>=<\gamma'',n>
\]

and this is how we define it.

\section*{Surfaces in Space}

Gauss bonnet theorem states that for a oriented surface, gaussian curvature $K$ is a function on a surface

\[
\int_M KdS=\int \frac{1}{a^2}dS=\frac{1}{a^2}4\pi a^2=4\pi
\]

where for the sphere $M$ of radius $a$ is in this case. \\

We define $\chi(M)$ as euler characteristic where for previous example it would be $2$.

\[
\int _M KdS=2\pi \chi(M)
\]

\section*{Directional Derivatives}

to compute directional derivatives, 

\[
x^i=p^i+ta^i
\]

we set the parametric equations for the line through $p$ in the direction $X_p$

let 

\[
a=(a^1,...,a^n)
\]

then direcitonal derivative $D_{X_p}f$ is 

\[
D_{X_p}f=\lim_{t\rightarrow0} \frac{f(p+ta)-f(p)}{t}=\frac{d}{dt}|_{t=0}f(p+ta)
\]
\[
=\Sigma\frac{\partial f}{\partial x^i}*\frac{dx^i}{dt}|=\Sigma\frac{\partial f}{\partial x^i}*a^i
\]
\[
=(\Sigma a^i\frac{\partial}{\partial x^i}|_p)f=X_pf
\]

and therefore directional derivative at $p$ of a $C^\infty$ vector field is

\[
D_{X_p}Y=\Sigma(X_pb^i)\frac{\partial}{\partial x^i}|_p
\]

The directional Derivative satisfies this property
\[
1. D_XY \text{ is} \; \mathcal{F}\text{-linear in } X \;\text{ and } \mathbb{R} \text{ linear in } Y
\]
\[
2. D_X(fY)=(Xf)Y+fD_XY
\]

we also define 

\[
[X,Y]=D_XY-D_YX
\]

and the quantity 

\[
T(X,Y)=D_XY-D_YX-[X,Y]
\]

as torsion. 

The definition for curvature is

\[
R(X,Y)=[D_X,D_Y]-D_{[X,Y]} =D_XD_Y-D_YD_X-D_{[X,Y]} \in End_\mathbb{R}(X(M))
\]

\section*{The Shape Operator}

for 

\[
X_p\in T_pM
\]

we define shape operator as

\[
L_p(X_p)=-D_{X_p}N
\]

\textbf{Lemma} let $M$ be the surface in $\mathbb{R}^3$ then 

\[
\forall X,Y\in \chi(M), <L(X),Y>=<D_XY,N>
\]

this is because

\[
0=X<Y,N>=<D_XY,N>+<Y,D_XN>=<D_XY,N>-<Y,L(X)>
\]

shape operator is self-adjoint as well, meaning

\[
<L(X),Y>=<X,L(Y)>
\]

we define second fundamental form as 

\[
\Pi(X,Y)=<L(X),Y>
\]

where first fundamental form is euclidean riemannian metric as we've defined earlier.

\section*{Affine connection}

affine connection is way of differentiating vector fields on manifold, where it is defined as

\[
\nabla_{X_p}f=X_pf
\]

and to do so, we define in the same way, nabla is defined as

\[
\nabla :\chi(M) \times \chi(M) \rightarrow \chi(M)
\]

then for

\[
\forall X,Y\in \chi(M), 
\]

it happens the same properties,

\[
\nabla_XY \text{ is } \mathcal{F}\text{-linear in } X
\]
\[
\nabla_X(fY)=(Xf)Y+f\nabla_XY
\]

it also has same definitions of torsion and curvature, where torsion is in $\chi(M)$ while euclidean connection on $\mathbb{R}^n$ so curvature is in endomorphism or the manifold

for this observation, affine connection $\nabla$ on a manifold $M$ gives linear map

\[
\chi(M)\rightarrow \text{End}_{\mathbb{R}}\chi(M)
\]

and this is map from

\[
X\rightarrow \chi(M)
\]

\section*{Theorema Egregium}

\[
i. K_p=<R_p(e_1,e_2)e_2,e_1>
\]
\[
ii. \text{gaussian curvature is an isometric invariant of smooth surfaces in }\mathbb{R}^3
\]

\section*{Generalizations to Hypersurfaces}

\textbf{Theorem:} let $M$ be a smooth hypersurface and $D$ the direcitonal derivative then

\[
\forall X,Y\in \chi(M), \nabla_XY=(D_XY)_{tan}
\]

where for torsion freeness, 

\[
D_XY=\nabla_XY+(D_XY)_{normal}
\]

and as we define forms of fundamental, first, second and third fundamental form is 

\[
I(X_p,Y_p)=<X_p,Y_p>
\]
\[
II(X_p,Y_p)=<L(X_P),Y_P>=<X_P,L(Y_P)>
\]
\[
III(X_P,Y_P)=<L^2(X_P),Y_P>=<L(X_P),L(Y_P)>=<X_P,L^2(Y_P)>
\]

\textbf{Prop} if $N$ is a smooth unit normal vector field on hypersurface $M$ in $\mathbb{R}^{n+1}$ and $X,Y\in\chi(M)$ then

\[
(D_XY)_{nor}=II(X,Y)N
\]

\section*{connections on vector bundles}

connection on a map is defined as

\[
\nabla : \chi(M) \times \Gamma(E)\rightarrow\Gamma(E)
\]

such that

\[
X\in \chi(M), \; s\in\Gamma(E)
\]

then states

\[
i. \; \nabla \text{ is } \mathcal{F}\text{-linear in } X \text{ and } \mathbb{R}\text{-linear in } s
\]
\[
ii. \; \nabla_X(fs)=(Xf)s+f\nabla_Xs
\]

for connection

\[
\nabla^U:\chi(U)\times \Gamma(U,E)\rightarrow\Gamma(U,E)
\]

it is defined as 

\[
\nabla_Xe_j=\sum w_j^i(X)e_i
\]

where

\[
R(X,Y)e_j=\sum \Omega_j^i(X,Y)e_i
\]

and 

\[
R(X,Y)=\nabla_X \nabla_Y-\nabla_Y\nabla_X-\nabla_{[X,Y]}
\]

and because of the definition of nabla, 

\[
R(X,Y)=(Xw_j^i(Y)-Yw_j^i(X)-w_j^i([X,Y]))e_i+(w_k^i(X)w_j^i(Y)-w_k^i(Y)w_j^k(X))e_i
\]
\[
=dw_j^i(X,Y)e_i+w_k^i\wedge w_j^k(X,Y)e_i
\]

therefore we can conclude that

\[
\Omega_j^i=dw_j^i+\sum_{k}w_k^i\wedge w_j^k
\]

also called as second structural equation, we also denote that $\tau$'s differential factor is $\omega$, therefore we can also conclude

first structural formula:

\[
\tau^i=d\theta^i+\sum_jw_j^i\wedge\theta^j
\]

\section*{Geodesics}

It is called \textbf{covariant derivative} if this is allowed

\[
\frac{D}{dt}: \Gamma(TM|_{c(t)})\rightarrow \Gamma(TM|_{c(t)})
\]

for this map, 

\[
\frac{DV}{dt} \text{ is } \mathbb{R} \text{ linear in } V
\]
\[
\frac{D(fV)}{dt}=\frac{df}{dt}V+f\frac{DV}{dt}
\]
\[
\frac{DV}{dt}(t)=\nabla_{c'(t)}V
\]

\textbf{Christoffel symbol} is a function that is of the connection $\nabla$ on the coordinate open set 

\[
(U,x^1,...,x^n)
\]

and defined as 

\[
\nabla_{\partial_i}\partial_j=\sum_{k=1}^n \Gamma_ij^k \partial_k 
\]

and for affine connection $\nabla$ on a manifold is torsion free iff every coordinate chart 

\[
(U,x^1,...,x^n) 
\]

the Christoffel symbol is symmetric in i and j:

\[
\Gamma_{ij}^k=\Gamma_{ji}^k
\]

ON a manifold with a connection, $c(t)$ the parametrized curve is a geodesic iff for

\[
(U,\phi)=(U,x^1,...,x^n)
\]

the components of 

\[
\ddot{y}^k+\sum_{i,j}\Gamma_{ij}^k\dot{y}^i\dot{y}^j=0
\]

this is because

\[
T(t)=c'(t)=\sum \dot{y}^j\partial_{j,c(t)}
\]
\[
\frac{DT}{dt}(t)=\sum_j\ddot{y}^j\partial_{j,c(t)}+\sum_j\dot{y}^j\nabla_{c'(t)}\partial_j
\]
\[
=\sum_k(\ddot{y}^k+\sum_{i,j}\dot{y}^i\dot{y}^j\Gamma_{ij}^k)\partial_k
\]

\textbf{Gauss-bonnet theorem for surface}

the euler characteristic is defined as

\[
\chi = V-E+F
\]

where we concluded

\[
\int_MKvol=2\pi \chi(M)
\]









\end{document}
